\section{Discussion}
\label{sec:discussion}

\subsection{What the agents attend to}

The completing deed is a textbook identity-dependent externality. Its holder is worse off if a particular rival gets it, so the holder's reservation price depends on who buys~\cite{jehiel1996nuclear}. The agents behave as if they know this. They state it in so many words (``\$450 is the price for the player who completes yellow''), and three kinds of behaviour follow it: recipients refuse proposer-completing openings more often than neutral ones, holders' asks rise for completing buyers in the pooled data, and cash concessions rarely buy such a deed. This is opponent attention of a particular kind. The agents reason about what a rival will be able to do with an asset, using holdings and cash that the prompt shows exactly. The records contain many such deductions about public capacity. They contain no comparable evidence about rivals' hidden beliefs or intentions, which the design does not measure.

The attention is mainly defensive. Holders price a rival's completion more consistently than buyers price their own. Recipients often refuse offers that would complete their own group when the counterparty also gains, and the refusals cite the rival's gain more often than any other consideration. In a game that only one player survives, relative position is what counts, so a refusal of a mutually completing swap is not irrational by itself. G15 shows the logic: Grok declines its own monopoly because Sol could build first. The records cannot tell us whether such refusals helped. A refusal followed by elimination, as in G15, is one path, not a counterfactual comparison.

\subsection{Attention without calibration}

The quantities the agents cite do not organise their responses. Rejection of proposer-completing openings is about 75\% whether the proposer would keep enough cash for a full first layer of houses or for one house only. Asking prices for completing deeds vary widely across seed blocks, and comparisons of one seller pricing one deed to two buyers are mixed. Counters reprice the same bundle whether or not completion is at stake. Several of the numbers inside the explanations are also wrong: \$520 does not buy nine yellow houses, a nine-square roll is not 5 in 36, and one house per green is not the minimum build.

A coherent reading is that the agents apply a categorical rule, roughly ``do not hand a rival a monopoly cheaply'', and then bargain over price with a generic routine. That reading matches the controlled two-party finding that LLM negotiators can model a counterparty accurately while their final prices follow opening anchors and local adjustment~\cite{cosentino2026counterparty}. Our setting is uncontrolled and multi-player, so it cannot test that mechanism. It does show the same separation between knowing what matters to the other side and pricing it consistently, in complete games with real consequences. We describe this as strategic attention. Strategic effectiveness would require a value reference for each decision, which this cohort lacks.

\subsection{A market that clears under distress}

If holders guard completing deeds and buyers do not pay more for them, few such deeds should trade except when a holder's circumstances change. That is what the settled sales show. Most completing deeds that sold for cash came from sellers holding less than \$100, and the six standard-condition sellers in that position still received a median of 2.25 times list. Buyers who complete a group this way often receive a mortgaged deed and must redeem it before building, which is where many of the unbuilt groups stall.

The collective consequence is that negotiation does not stop consolidation. Groups are completed by trade, purchase and inheritance; developed groups collect almost all rent; rent drives the debts that end players; and the first estate usually goes to a player who already leads. The agents see this chain and bargain over its last link. Grok sells ``insurance'' against Sol inheriting its yellows (G13), Opus withdraws a standing offer once the seller's estate is close (G18), and Gemini asks the table to stop funding a rival's houses (G20). The creditor rule makes a third player's bankruptcy a matter for everyone, and the agents' messages treat it that way.

\subsection{Conditions are part of the environment}

Opus is the one system whose bargaining does not look defensive, and it is also the system that played almost every decision under action-only recovery. We cannot separate the model from the condition. What matters for the other agents is that a willing trading partner existed. Opus opened 240 threads, sold half of the completing deeds that changed hands for cash, and sold 8 of the 14 distressed ones, seven of them below list. Grok's conflicting calls opened another 64 threads. Each participant's response condition therefore changed the offers that the others received. Analyses of multi-agent behaviour need to report such conditions as properties of the population, because deleting one participant's rows does not recover a clean experiment for the rest.

\subsection{The motivating questions}

The study began from eight capability questions and three broader ones (Appendix~\ref{app:questions}). The cohort answers the example in the first question most directly: agents do reject seemingly attractive trades because of what the trade gives an opponent, and their rejections and prices follow that consideration at the level of categories. The records also describe opponent attention to public capacity, bargaining tactics in context, liquidity management at every debt, and the concentration of assets through creditor succession. They do not measure calibrated beliefs about hidden information, negotiation efficiency, learning caused by feedback, the expected value of a price, transfer to unfamiliar opponents, or the relation between general and multi-agent capability. Those are unmeasured here, which is different from a finding that the agents lack them.

\section{Related work}
\label{sec:related}

\paragraph{Negotiation by language models.} NegotiationArena studies LLM bargaining in controlled two-party scenarios with tactic manipulations~\cite{bianchi2024negotiationarena}. Cosentino et al.\ give one side ranked information about the other's preferences in a multi-issue purchase negotiation and find that accurate counterparty modelling does not reliably improve the informed side's outcome, that final prices track opening anchors more than utility weights, and that sellers are more accommodating than buyers~\cite{cosentino2026counterparty}. Those designs have explicit utilities and can score agreements. Ours has none, but its bargaining happens inside a complete multi-player game in which the counterparty's gain affects third players and later turns.

\paragraph{Economic and board-game agents.} AgentMonopoly studies open-weight model families, sizes and personas in four-player Monopoly games with a 150-step limit and immediate elimination at negative cash, and argues that persona following should be judged by actions rather than language~\cite{wu2026agentmonopoly}. FinBoardBench evaluates financial decisions across board games, including Monopoly counteroffers, and its Monopoly bankruptcy treatment sells buildings to the bank~\cite{hu2026finboardbench}. DECE evaluates complete model-and-execution configurations across ten Eurogames with matched comparisons of execution frameworks~\cite{zhang2026dece}. Vending-Bench studies long-horizon operation of a simulated business~\cite{backlund2025vendingbench}. Its multi-agent communication study analyses 2,583 emails over twenty one-year simulations, checks factual claims against the simulator's event log, measures promise follow-through, and replicates its judging with other model families~\cite{li2026arena}. That study is the closest model for joining messages to state; our incidents are selected illustrations, not a population rate. The present paper adds a population analysis of bargaining over an identity-dependent externality, rebuilt from what each recipient saw, and follows the completed groups into development, debt and inheritance.

\paragraph{Externalities in exchange.} Jehiel, Moldovanu and Stacchetti model sales in which a non-acquirer suffers an externality that depends on who acquires the object, so that buyers' outside options are endogenous~\cite{jehiel1996nuclear}. We use the concept descriptively. Our agents are not mechanism designers, and Monopoly sellers are themselves exposed to the externality they price.

\paragraph{Reading agents' explanations.} Model explanations can be unfaithful to the computation that produced an answer~\cite{turpin2023unfaithful,lanham2023faithfulness}. We therefore treat private notes, public messages and provider-returned text as recorded channels, and draw behavioural conclusions from applied actions. Evaluations of complete agent systems also motivate reporting configuration and execution conditions with results~\cite{kapoor2026hal}.

\section{Limitations}
\label{sec:limits}

The cohort is small and structured. Five seeds, one fixed roster, a single cyclic seat order and visible model names limit what the patterns say about other opponents, prompts or seatings. Contrasts that hold in every seed block are still descriptions of these games.

The response conditions differ by system. Opus played most decisions without private notes or a corrective retry after a provider refusal, Grok's applied calls sometimes contradict its content, and all agents received some misattributed trade history. We report the main contrasts with and without the affected rows, but the conditions also shaped every other agent's opportunities, so no subset is a clean control.

The agents were not told the rules, and prices appear only in some prompts. Their valuations therefore mix strategic judgment with recall of the standard game, which the implementation sometimes departs from.

The comparisons are observational. Offer cells differ in deeds, prices, stages and states; the distress result conditions on settlement; asset value is a crude index. No value reference exists to say whether a refusal, a price or a bluff was good. The stated-reason codes come from two readers of one model family, one partly unblinded, and agreement between them is not validation. The eight incidents were selected to show mechanisms, and seven involve Grok, whose separate content and call channels make conflicts visible; they estimate no rate. The engine replay and the event tracker establish correspondence with the implemented rules, not with any other rule set.

\section{Conclusion}

In twenty complete games, four configured LLM agents bargained over the deed that would complete a rival's group as if its value depended on who bought it. They refused proposer-completing offers more often, rarely accepted cash concessions for such deeds, asked more from completing buyers in the pooled data, and said why. Their responses followed that category rather than the cash amounts they cited, and the deeds changed hands mainly when holders needed money. The groups then turned into income through building, and into control through creditor succession that usually rewarded an existing leader. The records establish what the agents attended to and what they did. Whether that defensive bargaining was good play is a question for a design with a value reference and varied opponents.
